\chapter{Manajemen Siklus Hidup API dan API Gateway}

\section*{Tujuan Pembelajaran}
\begin{itemize}
  \item Menjelaskan tahap siklus hidup API dan keputusan penting di tiap tahap.
  \item Menjelaskan fungsi \textit{API gateway} dan komponen platform manajemen API.
  \item Membandingkan alternatif platform manajemen API dengan kriteria yang dapat dipertanggungjawabkan.
\end{itemize}

\section{Pendahuluan}

% TODO
% Cakupan: perubahan API di satu rumah sakit dapat memutus integrasi mitra tanpa pemberitahuan,
% sementara tidak ada catatan mitra mana yang memakai versi mana.

\section{Siklus Hidup API}

% TODO
% Cakupan: tahap siklus hidup dari perancangan, pembangunan, publikasi, pemakaian, pemeliharaan,
% sampai pensiun; tingkat kematangan API menurut Medjaoui dkk.; keputusan dan pemilik keputusan di
% tiap tahap.
% Kasus Medikara: API jadwal dokter dari rancangan sampai dipakai klinik mitra.

\section{Versi dan Deprecation}

% TODO
% Cakupan: kebijakan versi (kapan versi baru diperlukan, cara penamaan versi); deprecation dan
% pensiun API; tenggat dan cara pemberitahuan kepada mitra; kaitan dengan kompatibilitas pada Bab 9.
% Kasus Medikara: mitra yang masih memakai API versi lama.

\section{API Gateway}

% TODO
% Cakupan: posisi API gateway di antara pengguna dan layanan; fungsi routing, autentikasi, rate
% limiting, kuota, transformasi, dan pencatatan; penerapan kebijakan keamanan dari Bab 6 di
% gateway; risiko gateway sebagai titik kegagalan tunggal.
% Kasus Medikara: satu gateway bersama untuk lima rumah sakit.
% Pendalaman teknis: menjalankan gateway open source di Docker dan menerapkan rate limiting serta
% validasi token pada satu API contoh.

\section{Platform Manajemen API}

% TODO
% Cakupan: komponen gateway, portal pengembang, analitik, dan pengelola kebijakan; katalog API;
% hubungan dengan portal pengembang pada Bab 11 dan metrik pada Bab 12.

\section{Memilih Platform Manajemen API}

% TODO
% Cakupan: pilihan open source, layanan cloud, dan hibrida; kriteria (fungsi, biaya, kendali,
% keterampilan tim, lokasi data, ketergantungan pada penyedia); matriks keputusan; ekonomi cloud
% dan FinOps dibahas rinci di IT40713.
% Kasus Medikara: pertimbangan lokasi data pasien dan kemampuan tim TI internal.

\section{API Gateway dan Service Mesh}

% TODO
% Cakupan: perbedaan lalu lintas dari luar ke dalam dan lalu lintas antarlayanan; peran gateway dan
% service mesh secara singkat; service mesh dan komunikasi antarlayanan dibahas rinci di IT40713.

\section{Aktivitas: Rancangan Manajemen API}

% TODO
% Tujuan aktivitas: merancang cara API target dikelola sepanjang hidupnya.
% Peran dalam tim: anggota berlatar bisnis dan hukum menyusun kebijakan versi dan tata cara
% pemberitahuan deprecation kepada mitra; anggota berlatar teknis menggambar posisi gateway dan
% aliran permintaan.
% Langkah: gambar arsitektur manajemen API target; tentukan kebijakan di gateway; susun kebijakan
% siklus hidup dan versi; bandingkan tiga alternatif platform manajemen API dalam matriks
% keputusan; tulis catatan keputusan.
% Luaran dan tonggak capstone: rancangan manajemen API dan kebijakan siklus hidup, masuk ke C3.
% Berkas sources yang direncanakan: sources/chapter10/kebijakan-siklus-hidup-api.md,
% sources/chapter10/matriks-pemilihan-platform-api.csv,
% sources/chapter10/docker-compose-gateway.yml (opsional)

\section{Ringkasan}

% TODO

\section{Latihan}

\begin{enumerate}
  \item Misalkan seorang mitra masih memakai API versi lama yang akan dipensiunkan. Susun langkah komunikasi dan tenggat yang adil bagi kedua pihak.
  \item Jelaskan fungsi \textit{API gateway} dengan analogi pos pemeriksaan, lalu sebutkan batas analogi tersebut.
  \item Bandingkan layanan manajemen API dari penyedia \textit{cloud} dengan produk \textit{open source} yang dipasang sendiri, dari sisi biaya, kendali, dan keterampilan tim.
\end{enumerate}
